\section{The model and the variational principle}\label{model:section}

Fix an even integer $p\ge2$ and $\alpha>0$. Let $\theta$ be a random real
function on $\{-1,1\}^p$ and let $h$ be a real random variable. Write
$B(\theta)=\max_s\abs{\theta(s)}$ and assume
\begin{equation}\label{model:integrability}
 \E B(\theta)+\E\abs h<\infty.
\end{equation}
For independent $M_N\sim\Pois(\alpha N)$, independent copies $\theta_k$ of
$\theta$, independent copies $h_i$ of $h$, and independent uniform indices
$i_{\ell,k}\in[N]=\{1,\ldots,N\}$, set
\begin{equation}\label{model:hamiltonian}
 \begin{gathered}
 -H_N(\sigma)=\sum_{k=1}^{M_N}
 \theta_k(\sigma_{i_{1,k}},\ldots,\sigma_{i_{p,k}})
 +\sum_{i=1}^N h_i\sigma_i,\\
 Z_N=\sum_{\sigma\in\{-1,1\}^N}e^{-H_N(\sigma)},
 \qquad F_N=\frac1N\E\log Z_N.
 \end{gathered}
\end{equation}
All families in this definition are mutually independent. In particular,
indices are sampled with replacement. The model parameters and disorder
laws are fixed as $N$ varies.

The interaction hypotheses are the existence of a representation
\begin{equation}\label{model:factorization}
 e^{\theta(s_1,\ldots,s_p)}
 =a\left(1+b\prod_{\ell=1}^p f_\ell(s_\ell)\right),
 \qquad
 \abs{b\prod_{\ell=1}^p f_\ell(s_\ell)}<1\quad\text{a.s.},
\end{equation}
for every spin tuple, where $a>0$ and $b$ are real random variables,
the $f_\ell:\{-1,1\}\to\R$ are independent copies of one random
function, $b$ is independent of them, and
\begin{equation}\label{model:positivity}
 \E[(-b)^n]\ge0\qquad(n\ge1).
\end{equation}
The expectations in \eqref{model:positivity} are interpreted as finite real
expectations. This is automatic in the nondegenerate case: independence
and the strict product bound force $b$ and $\max_s|f_1(s)|$ to be
essentially bounded, as shown in \cref{upper:bound}. If the product term
vanishes identically, the interaction is spin independent and one may
take $b=0$. The multiplier $a$ may depend on all the other variables in
the representation; no separate integrability of $\log a$ is assumed.

\subsection{The hierarchical cavity functional}

We give the power-mean form of the Panchenko--Talagrand functional
\cite[Section~5]{PT04}. A message $x$ specifies a probability distribution
on one spin. The site term inserts a new spin with its incident
interactions; the edge term corrects for interactions counted in that
insertion.

For $x\in\R$ let $q_x(s)=e^{xs}/(2\cosh x)$. Define
\begin{align}
 C_\theta(x_1,\ldots,x_p)
 &=\sum_{s\in\{-1,1\}^p}e^{\theta(s)}\prod_{\ell=1}^p q_{x_\ell}(s_\ell),
 \label{model:edge}\\
 U_\theta(x_1,\ldots,x_{p-1};\eps)
 &=\log\sum_{s\in\{-1,1\}^{p-1}}
 e^{\theta(s_1,\ldots,s_{p-1},\eps)}
 \prod_{\ell=1}^{p-1}q_{x_\ell}(s_\ell).
 \label{model:site-message}
\end{align}
Both $\abs{\log C_\theta}$ and $\abs{U_\theta}$ are bounded by $B(\theta)$,
uniformly in the messages.

Write $\mathcal P(E)$ for the Borel probability measures on $E$.
Let $\mathcal M_1=\mathcal P(\R)$ and
$\mathcal M_{j+1}=\mathcal P(\mathcal M_j)$, equipped with the topology of weak convergence and its Borel sigma-field. Fix $r\ge1$, $\zeta\in\mathcal M_{r+1}$, and
$0<m_1<\cdots<m_r<1$. For each message label $\omega$ independently, draw
\begin{equation}\label{model:chain}
 \eta_0^\omega\sim\zeta,\qquad
 \eta_{d+1}^\omega\mid\eta_d^\omega\sim\eta_d^\omega\ (0\le d\le r-2),
 \qquad x^\omega\mid\eta_{r-1}^\omega\sim\eta_{r-1}^\omega.
\end{equation}
At $r=1$, this means that each label first draws a random probability
measure $\eta_0$ with law $\zeta$, then draws $x$ with law $\eta_0$.
Higher levels repeat this sampling of measures. All these variables are
independent of the physical disorder. Use different
labels for $x_1,\ldots,x_p$ and for $x_\ell^j$, $j\ge1$, $1\le\ell<p$.
Take independent $k\sim\Pois(\alpha p)$, $\theta,\theta_1,\ldots$, and $h$,
and put
\begin{align}
 V_{\rm s}&=\frac12\sum_{\eps=\pm1}
 \exp\left(h\eps+\sum_{j=1}^k
 U_{\theta_j}(x_1^j,\ldots,x_{p-1}^j;\eps)\right),\label{model:vs}\\
 V_{\rm e}&=C_\theta(x_1,\ldots,x_p).\label{model:ve}
\end{align}
For $0\le d<r$, let $\mathscr F_d$ be generated by all the physical disorder in these
expressions and all $\eta_j^\omega$, $0\le j\le d$. Set
\begin{equation}\label{model:power-means}
 T_rV=V,\qquad
 T_dV=\left(\E[(T_{d+1}V)^{m_{d+1}}\mid\mathscr F_d]\right)^{1/m_{d+1}}
 \quad(d=r-1,\ldots,0).
\end{equation}
Only auxiliary messages are integrated in these conditional expectations.
For example, at depth one,
$T_0V=(\E_x[V^{m_1}\mid\boldsymbol\eta_0])^{1/m_1}$.
At depth two, each label draws
$\eta_0\sim\zeta$, then $\eta_1\mid\eta_0\sim\eta_0$, and then
$x\mid\eta_1\sim\eta_1$, and
\[
 T_0V=\left(\E_{\boldsymbol\eta_1}
   \left[\left(\E_{\boldsymbol x}
      [V^{m_2}\mid\boldsymbol\eta_1]\right)^{m_1/m_2}
      \,\middle|\,\boldsymbol\eta_0\right]\right)^{1/m_1}.
\]
Here bold symbols collect all message labels in $V$; physical disorder is
held fixed throughout, including in the two abbreviated displays.
Define
\begin{equation}\label{model:functional}
 \cB_r(\zeta,\mathbf m)
 =\log2+\E\log T_0V_{\rm s}
 -\alpha(p-1)\E\log T_0V_{\rm e},
 \qquad \Phi_r=\inf_{\zeta,\,\mathbf m}\cB_r(\zeta,\mathbf m).
\end{equation}
At $r=0$, the messages are independent with common law
$\zeta\in\mathcal P(\R)$ and the operators $T_d$ are omitted; this defines
$\cB_0(\zeta)$ and $\Phi_0$.

The quantities in \eqref{model:functional} are finite without any moment
assumption on the trial messages. Indeed, the power-mean operators preserve
pointwise multiplicative bounds, and
\begin{equation}\label{model:trial-bounds}
 \abs{\log T_0V_{\rm s}}\le\abs h+\sum_{j=1}^k B(\theta_j),
 \qquad \abs{\log T_0V_{\rm e}}\le B(\theta).
\end{equation}

\begin{theorem}\label{main:theorem}
Under \eqref{model:integrability}--\eqref{model:positivity}, the thermodynamic
limit exists and
\begin{equation}\label{main:formula}
 \lim_{N\to\infty}F_N=\inf_{r\ge0}\Phi_r.
\end{equation}
\end{theorem}

The right side is an infimum over finite hierarchies. The theorem does not
assume that a trial law attains it or that a limiting infinite-depth
Gibbs state has been constructed.
